\begin{table}[H]
\centering
\caption{Effect of the SDIL on Childhood Dental Decay: Continuous Treatment DiD}
\label{tab:main}
\begin{threeparttable}
\begin{tabular}{lccccc}
\toprule
 & (1) & (2) & (3) & (4) & (5) \\
\midrule
Post $\times$ IMD (std.) & -0.207 & -0.157 & -0.184 & -0.119 & -0.121 \\
 & (0.301) & (0.329) & (0.456) & (0.479) & (0.303) \\
\addlinespace
Post $\times$ IMD$^2$ & & -0.134 & & -0.137 & \\
 & & (0.245) & & (0.244) & \\
\addlinespace
Post $\times$ Obesity & & & Yes & Yes & \\
\midrule
LA FE & Yes & Yes & Yes & Yes & Yes \\
Year FE & Yes & Yes & Yes & Yes & Yes \\
Excl. COVID (2021/22) & & & & & Yes \\
Observations & 975 & 975 & 975 & 975 & 843 \\
R$^2$ (within) & 0.001 & 0.001 & 0.001 & 0.001 & 0.000 \\
Mean dep. var. & 27.2 & 27.2 & 27.2 & 27.2 & 27.5 \\
\bottomrule
\end{tabular}
\begin{tablenotes}[flushleft]
\small
\item \textit{Notes:} The dependent variable is the percentage of 5-year-olds with visually observable dental decay. Post equals one for survey waves 2018/19 onward. IMD~(std.) is the standardized Index of Multiple Deprivation 2019 score. Column~(2) adds a quadratic in IMD interacted with Post. Column~(3) controls for the pre-SDIL obesity rate interacted with Post. Column~(5) excludes the COVID-affected 2021/22 survey wave. Standard errors clustered at the local authority level in parentheses. \sym{*} $p<0.10$, \sym{**} $p<0.05$, \sym{***} $p<0.01$.
\end{tablenotes}
\end{threeparttable}
\end{table}

